% Generated draft: requires mathematical and source review.
\documentclass[11pt]{article}
\usepackage[margin=25mm]{geometry}
\usepackage{fontspec}
\setmainfont{Latin Modern Roman}
\usepackage{amsmath,amssymb,mathtools,unicode-math}
\setmathfont{Latin Modern Math}
\usepackage{enumitem,needspace,xurl,hyperref}
\setlength{\parindent}{0pt}
\setlength{\parskip}{5pt}
\setlength{\emergencystretch}{3em}
\newcommand{\N}{\mathbb{N}}
\newcommand{\Z}{\mathbb{Z}}
\newcommand{\Q}{\mathbb{Q}}
\newcommand{\R}{\mathbb{R}}
\newcommand{\C}{\mathbb{C}}
\newcommand{\F}{\mathbb{F}}
\newcommand{\A}{\mathbb{A}}
\newcommand{\PP}{\mathbb{P}}
\newcommand{\E}{\mathbb{E}}
\newcommand{\eps}{\varepsilon}
\newcommand{\dd}{\,\mathrm{d}}
\begin{document}
\section*{UnsolvedMath ID: 700811}
\textbf{Source problem code:} catalogue entry
\subsection{Definitions and mathematical statement}
Let \(\pi\) be a nonempty set of primes, and let \(\pi'\) denote its complement in the set of all primes. A finite group is a \(\pi\)-group if every prime dividing its order belongs to \(\pi\); the trivial group is included. Define a \(\pi'\)-group similarly.

For this draft, a finite group \(G\) is called \(\pi\)-soluble if it admits a finite normal series
\[
1=G_0\leq G_1\leq\cdots\leq G_n=G,
\qquad G_i\trianglelefteq G,
\]
in which each factor is either a soluble \(\pi\)-group or a \(\pi'\)-group. Its \(\pi\)-length, denoted \(\ell_\pi(G)\), is the minimum number of nontrivial \(\pi\)-group factors among all such series. In particular, \(\ell_\pi(1)=0\).

A Hall \(\pi\)-subgroup of \(G\) is a subgroup \(H\leq G\) such that \(H\) is a \(\pi\)-group and every prime dividing \([G:H]\) belongs to \(\pi'\). For a soluble group \(H\), put \(H^{(0)}=H\) and
\[
H^{(r+1)}=[H^{(r)},H^{(r)}]\qquad(r\geq0),
\]
and define its derived length by
\[
d(H)=\min\{r\geq0:H^{(r)}=1\}.
\]
Thus \(d(1)=0\). Under the convention for \(\pi\)-solubility above, Hall \(\pi\)-subgroups exist and are soluble.

\textbf{Question (L. A. Shemetkov, 1990; Kourovka Notebook Problem 11.119).} Is it true that, for every nonempty set of primes \(\pi\), every finite \(\pi\)-soluble group \(G\), and every Hall \(\pi\)-subgroup \(H\) of \(G\),
\[
\ell_\pi(G)\leq d(H)?
\]

\textit{Terminological qualification.} The supplied record does not define \(\pi\)-solubility or \(\pi\)-length. The definitions above use the convention requiring the \(\pi\)-factors to be soluble. Replacing this by the weaker requirement that the factors merely be \(\pi\)-groups or \(\pi'\)-groups can allow nonsoluble Hall \(\pi\)-subgroups, whose ordinary derived length is undefined. The intended convention should therefore be checked against the cited notebook.
\subsection{Short English statement}
Shemetkov asks whether, for every nonempty set of primes \(\pi\), the \(\pi\)-length of every finite \(\pi\)-soluble group is at most the derived length of its Hall \(\pi\)-subgroup. The draft makes the solubility convention explicit. The registry labels the problem open; this status is supplied rather than independently verified.
\subsection{Sources}
\begin{itemize}
\item Local UnsolvedMath registry, supplied complete source record: permanent ID \(700811\), exact source code \texttt{catalogue entry}, title \emph{Kourovka Notebook Problem 11.119}. Supplied status: open. Record updated on \texttt{2026-10-16T14:00:00Z}; upstream revision \texttt{d371ee917ac9647fbb06cd6fc529deec6f3dda82}.
\item Supplied source attribution: \emph{The Kourovka Notebook (Unsolved Problems in Group Theory)}, 11th issue (1990), Problem 11.119; proposed by L. A. Shemetkov.
\item The supplied background also identifies the 21st issue edition (2026), \texttt{arXiv:1401.0300v48}. The registry describes the collection as covering Issues 1--21 (1965--2026), edited by E. I. Khukhro and V. D. Mazurov. These bibliographic details are supplied attribution and have not been independently checked.
\end{itemize}
\end{document}
